\section{Experiments}
\label{sec:experiments}

\subsection{Datasets and Evaluation Metrics}
\label{sec:res-data}
We conduct experiments on three publicly available MSA benchmarks:
CMU-MOSI \citep{zadeh2017tensor}, CMU-MOSEI \citep{zadeh2018multi}, and
CH-SIMS v2 \citep{yu2021simsv2}.
CMU-MOSI and CMU-MOSEI provide English video clips annotated with continuous
sentiment scores in $[-3, 3]$, while CH-SIMS v2 is a Chinese benchmark that
emphasizes balanced modality-specific sentiment dominance and therefore serves
as a stronger test of sample-wise primary modality selection.
All three datasets provide pre-extracted language, acoustic, and visual
features following the standard preprocessing pipeline used in prior MSA
work \citep{mods2026align,yu2021selfmm}.

Following common MSA conventions, we report both \emph{regression-oriented}
scores---mean absolute error (MAE) and Pearson correlation (Corr) between
predicted and gold continuous labels---and \emph{classification-oriented}
scores derived from the same scalar predictions: binary accuracy (Acc-2)
after zero-thresholding, seven-class accuracy (Acc-7) on CMU-MOSI and
CMU-MOSEI, three- and five-class accuracy (Acc-3, Acc-5) on CH-SIMS v2, and
F1 (percent) where applicable.
MAE and Corr measure ordinal calibration on $[-3,3]$; Acc metrics summarize
hard decisions after thresholding or uniform binning, so they need not move
in lockstep with MAE when errors concentrate near class boundaries.
Higher values indicate better performance on every metric except MAE.
Unless stated otherwise, we select checkpoints by development MAE, but we
report the full regression--classification panel for each dataset.

\subsection{Implementation Details}
\label{sec:res-impl}
PRISM is implemented in PyTorch and trained with AdamW and linear warm-up.
We follow the two-stage curriculum described in \S\ref{sec:method-rib}:
stage~1 performs IB warm-up for a fixed number of epochs, after which stage~2
activates the routing supervision $\mathcal{L}_{\mathrm{route}}$.
The Gumbel--Softmax temperature is annealed linearly from
$\tau_{\mathrm{start}}$ to $\tau_{\mathrm{end}}$ across training, and after a
fixed burn-in epoch we additionally maintain exponential moving averages (EMA)
of the parameters for evaluation on the development and test splits.
Unless stated otherwise, we train and evaluate under the complete-modality
setting. For CH-SIMS v2 we additionally run an Optuna-based hyperparameter
search with the extended epoch range $n_{\mathrm{epochs}} \in [45, 85]$.
All experiments are conducted on a single RTX-4090-class GPU;
configuration files, random seeds, and trained checkpoints used to produce
the reported numbers will be released upon acceptance.

\subsection{Comparison with State-of-the-art Methods}
\label{sec:res-baseline}
Table~\ref{tab:lit} compares PRISM with representative MSA baselines on
CMU-MOSI and CMU-MOSEI, including the modality-invariant/-specific system
MISA \citep{hazarika2020misa}, the hierarchical mutual-information method
MMIM \citep{han2021mmim}, and a recent debiased IB system
CaMIB \citep{jiang2025towards}.
Numbers for the baselines are cited from the respective publications under
the complete-modality setting.

\paragraph{Regression.}
On CMU-MOSI, PRISM obtains the strongest MAE (0.605) among the listed systems,
with Corr (0.854) slightly below CaMIB (0.865), indicating strong ordinal
alignment that can diverge from a competitor's correlation--MAE trade-off on
the same clips.
On CMU-MOSEI, PRISM delivers the best Corr (0.814) while its MAE (0.599) trails
CaMIB's headline regression number; the gap is consistent with MOSEI's
noisier labels and with dev-MAE checkpointing, which does not explicitly
optimize correlation.

\paragraph{Classification.}
On CMU-MOSI, PRISM improves seven-way discretization accuracy (Acc-7 51.76)
relative to MISA/MMIM/CaMIB (46.6), reflecting finer-grained score bins, while
binary Acc-2/F1 remain within one point of CaMIB's best-published classifiers.
On CMU-MOSEI, PRISM leads Acc-2 (87.21) with competitive Acc-7 (47.60),
showing that the same configuration that helps hard thresholding also tracks
well on coarse multi-class cuts.

Overall, PRISM is strongest on the regression axis for MOSI and on
classification plus correlation for MOSEI, while staying close to prior art
elsewhere---consistent with primary-centric IB fusion improving calibration
and routing without sacrificing thresholded accuracy.
The results indicate that primary-centric IB-guided fusion is a viable
alternative to symmetric MulT-style cross-modal attention
\citep{tsai2019multimodal} and to the full causal modeling paradigm explored
by CaMIB \citep{jiang2025towards}.

\begin{table*}[t]
  \centering
  \footnotesize
  \setlength{\tabcolsep}{3pt}
  \begin{tabular}{lccccccccc}
    \toprule
    & \multicolumn{5}{c}{\textbf{CMU-MOSI}} & \multicolumn{4}{c}{\textbf{CMU-MOSEI}} \\
    \cmidrule(lr){2-6} \cmidrule(lr){7-10}
    \textbf{Model} & Acc-2 & Acc-7 & F1 & MAE & Corr & Acc-2 & Acc-7 & MAE & Corr \\
    \midrule
    MISA~\citep{hazarika2020misa}$^{\dagger}$
      & 83.4 & 42.3 & 83.6 & 0.855 & 0.761
      & 85.5 & 52.2 & 0.555 & 0.756 \\
    MMIM~\citep{han2021mmim}$^{\dagger}$
      & 86.0 & 46.6 & 85.9 & 0.700 & 0.800
      & 85.9 & 54.2 & 0.526 & 0.772 \\
    CaMIB~\citep{jiang2025towards}$^{\dagger}$
      & \textbf{88.5} & 46.6 & \textbf{88.3} & 0.617 & \textbf{0.865}
      & 86.5 & \textbf{55.0} & \textbf{0.525} & 0.784 \\
    \midrule
    \textbf{PRISM (Ours)}
      & 87.94 & \textbf{51.76} & 87.90 & \textbf{0.605} & 0.854
      & \textbf{87.21} & 47.60 & 0.599 & \textbf{0.814} \\
    \bottomrule
  \end{tabular}
  \caption{Comparison on CMU-MOSI and CMU-MOSEI (complete-modality): regression
    columns (MAE, Corr) and classification columns (Acc-2, Acc-7, F1).
    $\dagger$: cited from originals. Bold: best per column.}
  \label{tab:lit}
\end{table*}

\subsection{CH-SIMS v2 and Hyperparameter Search}
\label{sec:res-sims}
Because CH-SIMS v2 exposes modality-balanced sentiment dominance, it is a
stronger test of the primary-modality-centric design than the
language-dominant CMU-MOSI and CMU-MOSEI.
We therefore perform an Optuna-based hyperparameter search with tiered search
spaces, still using development MAE as the selection objective but reporting
the full classification panel on the held-out split.
Table~\ref{tab:hpo} lists the global bounds of the tier-1 search.
The tier-1 search (80 trials) converges to a best configuration with
development MAE of 0.3192, test Acc-2 of 78.05\%, Acc-3 of 72.92\%, and test
Acc-5 of 52.80\%;
a subsequent tier-2 search further widens the bounds on $\alpha_{\mathrm{IB}}$,
weight decay, warm-up, and EMA start.

\begin{table}[t]
  \centering
  \scriptsize
  \begin{tabular}{ll}
    \toprule
    \textbf{Group} & \textbf{CH-SIMS v2 tier-1 search bounds} \\
    \midrule
    Epochs & $n_{\mathrm{epochs}} \in [45, 85]$ (integer) \\
    Batch configs & index over $\{(8,4),(16,2),(16,4),(32,1),(32,2),(64,1)\}$ \\
    Learning rates &
      $\mathrm{LR} \in [5{\times}10^{-6}, 5{\times}10^{-5}]$ (log); \\
      & $\mathrm{LR}_{\mathrm{IG}} \in [5{\times}10^{-5}, 2{\times}10^{-3}]$ (log) \\
    IB / architecture &
      $\beta_{\mathrm{IB}} \in [4, 64]$ (log); \\
      & layers $\in \{2, 3, 4, 5\}$;
        bottleneck $\in \{64, 96, 128, 192\}$ \\
    Regularization &
      $\rho_{\mathrm{mse}} \in [0, 2]$; dropout $\in [0.05, 0.40]$ \\
    \bottomrule
  \end{tabular}
  \caption{Tier-1 Optuna search bounds on CH-SIMS v2.}
  \label{tab:hpo}
\end{table}

Table~\ref{tab:sims} compares PRISM under a default hyperparameter
configuration with the Optuna-selected configuration.
The search yields a clear \emph{regression} gain (MAE 0.3800 $\rightarrow$
0.3192) together with stronger \emph{multi-way classification} (Acc-3/Acc-5),
while Acc-2 dips slightly---illustrating the MAE--Acc trade-off when scores
hover near the zero threshold under dev-MAE selection.

\begin{table}[t]
  \centering
  \footnotesize
  \begin{tabular}{lcccc}
    \toprule
    \textbf{Configuration} & Acc-2 & Acc-3 & Acc-5 & MAE $\downarrow$ \\
    \midrule
    Default configuration    & \textbf{79.40} & ---   & 46.20 & 0.3800 \\
    Optuna-selected (tier-1) & 78.05 & \textbf{72.92} & \textbf{52.80} & \textbf{0.3192} \\
    \bottomrule
  \end{tabular}
  \caption{CH-SIMS v2: PRISM under a default hyperparameter configuration
    vs.\ the Optuna-selected tier-1 configuration (dev-selected). Bold
    indicates the better value per column.}
  \label{tab:sims}
\end{table}

\subsection{Ablation Study}
\label{sec:res-abl}
To verify the contribution of individual components, we conduct an ablation
on CMU-MOSI by removing one design element at a time.
Table~\ref{tab:abla} lists MAE under the same dev-MAE selection protocol as
the full model; we omit Acc columns for brevity.
We discuss each variant below.

\begin{table}[t]
  \centering
  \footnotesize
  \begin{tabular}{lc}
    \toprule
    \textbf{Setting on CMU-MOSI} & \textbf{MAE} $\downarrow$ \\
    \midrule
    Full PRISM & \textbf{0.605} \\
    \midrule
    (i) w/o EMA at evaluation & --- \\
    (ii) w/o stage-2 $\mathcal{L}_{\mathrm{route}}$ & --- \\
    (iii) Fixed text-primary (no MSelector) & --- \\
    \bottomrule
  \end{tabular}
  \caption{Ablation of PRISM components on CMU-MOSI.
    Bold indicates the reference full-model score.}
  \label{tab:abla}
\end{table}

\noindent \textbf{Effect of EMA at evaluation}.
Removing EMA reverts evaluation to the last-epoch parameters and exposes
epoch-to-epoch oscillation of development MAE around the stage-1/stage-2
transition, which in turn destabilizes checkpoint selection and degrades
dev-selected test MAE.

\noindent \textbf{Effect of stage-2 routing supervision}.
Without the routing supervision $\mathcal{L}_{\mathrm{route}}$, the MSelector
receives gradients only through the primary path of the fusion stack.
In this regime the selector tends to collapse toward the modality whose
gradient dominates---typically the language modality on CMU-MOSI---thereby
reducing the model to a near-static language-primary variant and negating
the sample-wise dynamics that the MSelector is meant to capture.

\noindent \textbf{Effect of MSelector}.
Replacing the learned MSelector with a hard language-primary assignment
disables sample-wise modality selection entirely.
On clips where acoustic or visual cues carry the decisive affective signal,
the fixed-text variant must rely on auxiliary cross-attention alone to
override the language prior, which is insufficient when lexical content is
neutral or polarity-conflicting with the non-language modalities.

\subsection{Case Study and Further Analysis}
\label{sec:res-qual}

\noindent \textbf{Qualitative failure modes}.
On CMU-MOSI, a substantial fraction of high-MAE clips mix lexical sentiment
with vocal sarcasm.
In these cases the MSelector correctly shifts weight toward the acoustic
modality, but the final prediction can still disagree with the label when
the three modalities carry conflicting polarities.
Such polarity-conflicting clips are analogous to the hard cases previously
analyzed in hierarchical IB fusion \citep{ithp2024align}.

\noindent \textbf{CMU-MOSI vs.\ CMU-MOSEI}.
The larger scale of CMU-MOSEI stabilizes batch-level routing statistics but
also extends the wall-clock training time.
We observe no stage-transition forgetting when the Gumbel temperature stays
above $0.5$ for the first half of training; annealing too aggressively tends
to collapse the routing distribution to a nearly static modality choice,
echoing the need for carefully staged objectives also reported by
token-level IB methods \citep{cyin2025align}.

\noindent \textbf{Regression vs.\ classification metrics.}
We select checkpoints by development-set MAE by default, following prior MSA
practice \citep{hazarika2020misa,han2021mmim}, because it directly reflects
ordinal error on $[-3,3]$.
When deployments care more about hard decisions (Acc-2) or fine bins
(Acc-7/Acc-5), a composite dev score that mixes MAE with one or more Acc
terms, or a separate dev run with Acc-tuned selection, reduces the known
mismatch between optimal regressors and optimal classifiers on discretized
outputs.

\noindent \textbf{Routing statistics}.
We monitor the batch-averaged selector masses $(w_t, w_a, w_v)$ throughout
training to detect potential routing collapse and to diagnose the dynamics of
sample-wise primary modality choice.
